% Gliederung Stand 30.10.2026. Grundlagen, ~6 Seiten.
\chapter{Humanoide Roboter und der Unitree G1}
\label{chap:humanoide}
% Kurze Kapiteleinleitung vor 2.1.

\section{Humanoide Roboter in der Produktion}
\label{sec:humanoide_produktion}

\subsection{Aufbau und Freiheitsgrade humanoider Roboter}
\label{sec:humanoide_aufbau}
% Kandidaten: Saeedvand2019 (Körperteile, Roll/Nick/Gier, mind. 20 DoF, DoF vs. Kosten), Westervelt2018 (Begriffe Zweibeiner, voll-/unteraktuiert).

\subsection{Einsatzszenarien und Hemmnisse bei der Einführung}
\label{sec:humanoide_einsatz}
% Kosten, Risiko, fehlendes Know-how. Kandidaten: Tong2024 (Industrie-Einsatz, Kosten, Sicherheit), Saeedvand2019 (Anwendungen, Aktorik als Grenze), Yu2026 (G1 in Unterhaltung/Bühnentanz), Westervelt2018 (Beine in menschengemachten Umgebungen; Stand 2007: kaum Nutzen über Neuheit hinaus).

\section{Hardware des Unitree G1}
\label{sec:g1_hardware}
% 29 Gelenke, Aktorik, Sensorik (IMU, Lidar), Inspire-FTP-Hände mit Taktilsensorik, Rechner. Kandidat: Saeedvand2019 (IMU als wichtigster Sensor).

\section{Software-Schnittstellen des Unitree G1}
\label{sec:g1_software}

\subsection{Unitree SDK2 und DDS-Kommunikation}
\label{sec:g1_sdk}
% LowCmd/LowState.

\subsection{Onboard-Regler und Zustandsmaschine}
\label{sec:g1_onboard}
% High-Level-API (LocoClient), FSM-Zustände.

\subsection{Oberkörper-Schnittstelle arm\_sdk}
\label{sec:g1_armsdk}
% Gewichtsslot zum Überblenden zwischen Onboard-Regler und eigener Armregelung.
